\section*{Hoofdstuk \ref{ch:g8:speed} --- Evenredigheid, snelheid en gemiddelden}

\begin{solution}{exo:g8:speed:1}
$\frac{150}{2} = 75$ km/h. \quad
$45$ min $= 0.75$ h: $\frac{27}{0.75} = 36$ km/h. \quad
$\frac{100}{12.5} = 8$ m/s.
\end{solution}

\begin{solution}{exo:g8:speed:2}
$2 \times 85 = 170$ km. \quad
$40$ min $= \frac23$ h: $90 \times \frac23 = 60$ km. \quad
$\frac{315}{90} = 3.5$ h $= 3$ h $30$ min.
\end{solution}

\begin{solution}{exo:g8:speed:3}
$72 \div 3.6 = 20$ m/s. \quad
$5 \times 3.6 = 18$ km/h. \quad
$108 \div 3.6 = 30$ m/s.
\end{solution}

\begin{solution}{exo:g8:speed:4}
$600 \div 15 = 40$ min. \quad
$2$ h $30 = 150$ min: $15 \times 150 = 2250$ L.
\end{solution}

\begin{solution}{exo:g8:speed:5}
$\frac{3.30}{1.2} = 2.75$ euro/kg tegenover $\frac{2.16}{0.8} = 2.70$ euro/kg:
het tweede aanbod is (net) de betere koop.
\end{solution}

\begin{solution}{exo:g8:speed:6}
\[
\bar v = \frac{25 \times 12.4 + 35 \times 10}{60}
= \frac{310 + 350}{60} = \frac{660}{60} = 11 .
\]
\end{solution}

\begin{solution}{exo:g8:speed:7}
$d = 340 \times 6 = 2040$ m $\approx 2$ km.
\end{solution}

\begin{solution}{exo:g8:speed:8}
Afstanden: $2 \times 5 = 10$ km, en daarna $1.5 \times 4 = 6$ km: samen $16$ km.
Verstreken tijd: $2 + 0.5 + 1.5 = 4$ h. \omterm{def:g8:speed:def}{Gemiddelde snelheid}:
$\frac{16}{4} = 4$ km/h.
\end{solution}

\begin{solution}{exo:g8:speed:9}
Gewogen:
$\frac{3 \times 15 + 2 \times 9 + 1 \times 12}{3 + 2 + 1}
= \frac{45 + 18 + 12}{6} = \frac{75}{6} = 12.5$.
Ongewogen: $\frac{15 + 9 + 12}{3} = 12$. Ze verschillen doordat de
wegingsfactoren het cijfer $15$ drie keer zo zwaar laten wegen als het cijfer
$12$.
\end{solution}

\begin{solution}{exo:g8:speed:10}
\emph{1.} Tijden: $\frac{120}{80} = 1.5$ h en $\frac{120}{120} = 1$ h. \omterm{def:g8:speed:def}{Gemiddelde
snelheid}: $\frac{240}{2.5} = 96$ km/h.

\emph{2.} De auto brengt \emph{meer tijd} door bij $80$ km/h ($1.5$ h) dan bij
$120$ km/h ($1$ h): de langzame \omterm{def:g8:speed:def}{snelheid} weegt zwaarder in het met tijd \omterm{def:g8:speed:average}{gewogen
gemiddelde}, en trekt het onder het midden $100$.
\end{solution}

\begin{solution}{exo:g8:speed:11}
De kloof van $18$ km sluit met $5 + 4 = 9$ km/h: ze ontmoeten elkaar na
$\frac{18}{9} = 2$ h, om $12{:}00$. Anna heeft dan $2 \times 5 = 10$ km gelopen:
ze ontmoeten elkaar $10$ km van het dorp van Anna (en $8$ km van dat van Boris ---
controle: $10 + 8 = 18$).
\end{solution}

\begin{solution}{pb:g8:speed:1}
\textbf{1.} Volgens \cref{def:g8:speed:def} is $t_1 = \frac{d}{v_1}$ en
$t_2 = \frac{d}{v_2}$; de hele tocht duurt $\frac{d}{v_1} + \frac{d}{v_2}$ voor
een afstand van $2d$.

\textbf{2.} Breng $d$ in de \omterm{def:g4:fractions:def}{noemer} buiten haakjes en streep het weg:
\[
\bar v
= \frac{2d}{d\left(\frac{1}{v_1} + \frac{1}{v_2}\right)}
= \frac{2}{\ \frac{1}{v_1} + \frac{1}{v_2}\ } .
\]
Binnenin op gemeenschappelijke \omterm{def:g4:fractions:def}{noemer}:
$\frac{1}{v_1} + \frac{1}{v_2} = \frac{v_2 + v_1}{v_1 v_2}$, en delen is
vermenigvuldigen met het omgekeerde (\cref{thm:g8:fractions:div}):
\[
\bar v = 2 \times \frac{v_1 v_2}{v_1 + v_2}
= \frac{2\, v_1 v_2}{v_1 + v_2} .
\]

\textbf{3.} $\dfrac{2 \times 30 \times 15}{30 + 15} = \dfrac{900}{45} = 20$ km/h,
zoals in \cref{ex:g8:speed:twolegs}; en
$\dfrac{2 \times 80 \times 120}{80 + 120} = \dfrac{19\,200}{200} = 96$ km/h,
zoals gevonden in \cref{exo:g8:speed:10}.

\textbf{4.} Het gemiddelde hangt alleen van de twee snelheden af, niet van $d$:
een heen-en-terugrit van $1$ km en een van $100$ km met dezelfde twee snelheden
hebben precies dezelfde \omterm{def:g8:speed:def}{gemiddelde snelheid}. (Daarom kon de formule op tochten van
verschillende lengte nagegaan worden.)

\textbf{5.} Bij een tijd $T$ met elke \omterm{def:g8:speed:def}{snelheid} zijn de afstanden $v_1 T$ en
$v_2 T$, dus is
\[
\bar v = \frac{v_1 T + v_2 T}{2T}
= \frac{(v_1 + v_2)\,T}{2T}
= \frac{v_1 + v_2}{2} .
\]
Een gemiddelde van snelheden wordt altijd gewogen met de \emph{tijd}
(\cref{def:g8:speed:average}). Bij gelijke tijden dragen de twee snelheden
gelijke gewichten: het gewone midden. Bij gelijke afstanden neemt de langzamere
\omterm{def:g8:speed:def}{snelheid} \emph{meer} tijd in, en dus meer gewicht --- en schuift het gemiddelde
naar haar toe.

\textbf{6.} Voor $(30, 15)$: $H = 20$ en $A = 22.5$. Voor $(80, 120)$: $H = 96$
en $A = 100$. Het rekenkundig gemiddelde wint beide keren.

\textbf{7.} Op de gemeenschappelijke \omterm{def:g4:fractions:def}{noemer} $2(v_1 + v_2)$:
\[
A - H
= \frac{(v_1 + v_2)^2 - 4\, v_1 v_2}{2\,(v_1 + v_2)} ,
\]
en de \omterm{def:g4:fractions:def}{teller} werk je uit met de merkwaardige \omterm{def:g2:mult:def}{producten} van
\cref{pb:g8:equations:1}:
\[
(v_1 + v_2)^2 - 4 v_1 v_2
= v_1^2 + 2 v_1 v_2 + v_2^2 - 4 v_1 v_2
= v_1^2 - 2 v_1 v_2 + v_2^2
= (v_1 - v_2)^2 .
\]

\textbf{8.} Een kwadraat is nooit negatief, en $2(v_1 + v_2)$ is positief: dus is
$A - H \geq 0$, dat wil zeggen $H \leq A$, voor elk paar positieve snelheden.
Gelijkheid vraagt $(v_1 - v_2)^2 = 0$, dat wil zeggen $v_1 = v_2$: zodra de twee
snelheden verschillen, zakt het gemiddelde van de heen-en-terugrit strikt onder
het midden --- de verrassing is een stelling.

\textbf{9.} Vermenigvuldigen:
\[
H \times A
= \frac{2\, v_1 v_2}{v_1 + v_2} \times \frac{v_1 + v_2}{2}
= v_1 v_2 ,
\]
na het wegstrepen van de factor $2$ en de factor $v_1 + v_2$.

\textbf{10.} Met de formule: $\bar v = \dfrac{2 \times 40 \times 60}{40 + 60}
= \dfrac{4800}{100} = 48$ km/h. Langs de lange weg: $\frac{60}{40} = 1.5$ h en
$\frac{60}{60} = 1$ h, dus $120$ km in $2.5$ h: $\frac{120}{2.5} = 48$ km/h. Ze
stemmen overeen.

\textbf{11.} Gemiddeld $30$ km/h over $30$ km betekent een totale tijd van ten
hoogste $\frac{30}{30} = 1$ h. De eerste $15$ km bij $15$ km/h kosten al
$\frac{15}{15} = 1$ h: het \emph{hele} tijdbudget is op, met nog de \omterm{def:g3:division:half}{helft} van het
parcours te gaan.

\textbf{12.} Bij $90$ km/h kost de tweede \omterm{def:g3:division:half}{helft} $\frac{15}{90} = \frac16$ h, dus
duurt de wedstrijd $\frac76$ h en is
\[
\bar v = \frac{30}{\ 7/6\ } = \frac{180}{7} \approx 25.7
\text{ km/h} < 30 .
\]
Wat de \omterm{def:g8:speed:def}{snelheid} op de tweede \omterm{def:g3:division:half}{helft} ook is, haar duur is een positief getal, dus
komt de totale tijd boven $1$ h en blijft het gemiddelde strikt onder $30$ km/h.
(Met de formule: $\frac{2 \times 15 \times v_2}{15 + v_2} = 30$ zou
$30 v_2 = 450 + 30 v_2$ opleggen, dat wil zeggen $0 = 450$ --- geen enkele
\omterm{def:g8:speed:def}{snelheid} $v_2$ werkt.)

\textbf{13.} Een gemiddelde van $20$ km/h laat $\frac{30}{20} = 1.5$ h toe; na het
eerste uur blijven er $15$ km over voor $0.5$ h: ze heeft
$\frac{15}{0.5} = 30$ km/h nodig. Controle:
$H(15, 30) = \frac{2 \times 15 \times 30}{45} = \frac{900}{45} = 20$ km/h.

\textbf{14.} Per minuut vullen de kranen $\frac{1}{20}$ en $\frac{1}{30}$ van de
kuip; samen
\[
\frac{1}{20} + \frac{1}{30}
= \frac{3}{60} + \frac{2}{60}
= \frac{5}{60} = \frac{1}{12}
\]
van de kuip per minuut: de kuip is in $12$ min vol. En
$H(20, 30) = \frac{2 \times 600}{50} = 24 = 2 \times 12$: de twee kranen samen
doen er precies \emph{de \omterm{def:g3:division:half}{helft} van het \omterm{pb:g8:speed:1}{harmonisch gemiddelde}} van hun soloduren
over --- debieten tellen op, dus zit je weer op het terrein van het \omterm{pb:g8:speed:1}{harmonisch
gemiddelde}.

\textbf{15.} Dezelfde afstand in beide richtingen, dus is
\[
\bar v = \frac{2 \times 900 \times 700}{900 + 700}
= \frac{1\,260\,000}{1\,600} = 787.5 \text{ km/h},
\]
onder de $800$ km/h van windstil weer. In één zin: het vliegtuig vliegt
\emph{langer} tegen de wind in dan met de wind mee, dus krijgt het langzame stuk
het grotere gewicht in het met tijd \omterm{def:g8:speed:average}{gewogen gemiddelde} --- wind kan een
heen-en-terugvlucht alleen maar verlengen.
\end{solution}
